\labsection{2}{Files, metadata, and reliable CLI utilities}{sec:lab02}
\begin{labproject}{Read file metadata, copy bytes safely, and report search status}

\textbf{Coverage.} Chapter~\ref{ch:files-cli}: contents vs. metadata, file descriptors, exit status.\\
\textbf{Goal.} Run three small utilities and verify them with standard Linux commands.\\
\textbf{Explain.} Explain what each utility promises to the user and how it reports failure.

\section*{Source code}
\labsourcecard{\texttt{file\_info.c}}{Read metadata with \texttt{stat()}.}{The program does not read the file's contents; it reads information \emph{about} the file.}
\lstinputlisting[style=utpcnumbered]{all_experiments/Lab02/file_info.c}

\labsourcecard{\texttt{file\_copy.c}}{Copy bytes using POSIX file descriptors.}{The outer loop reads chunks; the inner loop handles a write that may complete only part of a chunk.}
\lstinputlisting[style=utpcnumbered]{all_experiments/Lab02/file_copy.c}

\labsourcecard{\texttt{mini\_grep.c}}{Search a text file and return a useful exit status.}{\texttt{0}=match, \texttt{1}=no match, \texttt{2}=usage/file error.}
\lstinputlisting[style=utpcnumbered]{all_experiments/Lab02/mini_grep.c}

\section*{Task 1 --- Build and inspect metadata}
\begin{lstlisting}[style=utpbash]
cd all_experiments/Lab02
gcc -Wall -Wextra -std=c11 file_info.c -o file_info
gcc -Wall -Wextra -std=c11 file_copy.c -o file_copy
gcc -Wall -Wextra -std=c11 mini_grep.c -o mini_grep
./file_info sample.txt
stat sample.txt
\end{lstlisting}
Match at least three fields between the program and \texttt{stat}: size, inode, link count, or permissions.

\section*{Task 2 --- Verify the copy, not just the message}
\begin{lstlisting}[style=utpbash]
./file_copy sample.txt copy.txt
cmp sample.txt copy.txt
echo $?
ls -l sample.txt copy.txt
\end{lstlisting}
A \texttt{cmp} status of 0 is independent evidence that the byte contents are equal.

\section*{Task 3 --- Treat exit status as part of the interface}
\begin{lstlisting}[style=utpbash]
./mini_grep kernel sample.txt; echo "status=$?"
./mini_grep pineapple sample.txt; echo "status=$?"
./mini_grep kernel missing.txt; echo "status=$?"
\end{lstlisting}
Create a three-row table: command, visible output, exit status, meaning.

\section*{Task 4 --- Controlled edit in Vim}
Open \texttt{mini\_grep.c}. Change the match output from \texttt{FILE:LINE: text} to begin with \texttt{MATCH }. Save, recompile, and confirm that the exit status rules did not change.

\begin{keyidea}
A professional CLI tool has two audiences: a person reads its output; another program can read its exit status.
\end{keyidea}

\begin{quickcheck}
\begin{enumerate}
  \item Which information belongs to file metadata rather than file contents?
  \item Why must a copy program check the return value from \texttt{read()} and \texttt{write()}?
  \item Why is exit status useful even when a command already prints a message?
\end{enumerate}
\end{quickcheck}

\begin{verification}
\begin{itemize}
  \item At least three metadata fields from \texttt{file\_info} match \texttt{stat}.
  \item \texttt{cmp} returns status~0 for the original and the copy.
  \item \texttt{mini\_grep} returns 0 for a match, 1 for no match, and 2 for a missing file.
  \item After the Vim edit the output begins with \texttt{MATCH} and the exit-status rules are unchanged.
\end{itemize}
\end{verification}

\begin{labdeliverable}
Submit: metadata comparison, copy verification, the three-case exit-status table, and one screenshot or terminal capture showing the modified \texttt{MATCH} output.
\end{labdeliverable}
\end{labproject}
